\Article{21}
Since the observations have supplied the equations (approximately true) \(v=0,v'=0,v''=0\) etc., to derive the value of the unknown \(x\) from them we must use a combination of those equations of the form
\[
\kappa v+\kappa'v'+\kappa''v''+\text{etc.}=0
\]
which gives \(x\) the coefficient \(1\) and eliminates the remaining unknowns \(y,z\) etc.; by art. 18 this determination must be assigned the weight
\[
=\frac1{\kappa\kappa+\kappa'\kappa'+\kappa''\kappa''+\text{etc.}}
\]
It therefore follows from the preceding article that the most suitable determination will be the one in which we put \(\kappa=\alpha,\kappa'=\alpha',\kappa''=\alpha''\) etc. In this way \(x\) receives the value \(A\), and evidently the same value can also (without knowing the multipliers \(\alpha,\alpha',\alpha''\) etc.) be derived directly by elimination from the equations \(\xi=0,\eta=0,\zeta=0\) etc. The weight to be assigned to this determination will be \(=\frac1{[\alpha\alpha]}\), or the mean error to be feared in it
\[
=m\BareRoot{p[\alpha\alpha]}=m'\BareRoot{p'[\alpha\alpha]}=m''\BareRoot{p''[\alpha\alpha]}\text{ etc.}
\]

In precisely the same way, the most suitable determination of the remaining unknowns \(y,z\) etc. will give them the same values as result from elimination from the same equations \(\xi=0,\eta=0,\zeta=0\) etc.

Denoting the indefinite sum \(vv+v'v'+v''v''\text{ etc.}\), or, what amounts to the same thing, this sum
\[
p(V-L)^2+p'(V'-L')^2+p''(V''-L'')^2+\text{etc.}
\]
\SourceLeaf{030}by \(\Omega\), we see that \(2\xi,2\eta,2\zeta\) etc. are the partial differential quotients of the function \(\Omega\), namely
\[
2\xi=\frac{\mathrm{d}\Omega}{\mathrm{d}x},\quad2\eta=\frac{\mathrm{d}\Omega}{\mathrm{d}y},\quad2\zeta=\frac{\mathrm{d}\Omega}{\mathrm{d}z}\text{ etc.}
\]
The values of the unknowns resulting from the most suitable combination of observations, which we may conveniently call the \textit{most plausible values}, will therefore be identical with those for which \(\Omega\) attains its minimum value. Now \(V-L\) expresses indefinitely the difference between the computed and observed value. Thus the most plausible values of the unknowns will be those which make the sum of the squares of the differences between the observed and computed values of the quantities \(V,V',V''\) etc., multiplied by the weights of the observations, a minimum: a principle that we had established by a very different route in the \textit{Theoria Motus Corporum Coelestium}. And if the relative precision of the individual determinations is also to be assigned, we must, by indefinite elimination from equations (III), express \(x,y,z\) etc. themselves in the following form:
\[
\left.\begin{aligned}
x&=A+[\alpha\alpha]\xi+[\alpha\beta]\eta+[\alpha\gamma]\zeta+\text{etc.}\\
y&=B+[\beta\alpha]\xi+[\beta\beta]\eta+[\beta\gamma]\zeta+\text{etc.}\\
z&=C+[\gamma\alpha]\xi+[\gamma\beta]\eta+[\gamma\gamma]\zeta+\text{etc.}
\end{aligned}\right\}\quad\text{(VII)}
\]
\begin{center}etc.\end{center}
In this way the most plausible values of the unknowns \(x,y,z\) etc. will be respectively \(A,B,C\) etc., and the weights to be assigned to these determinations \(\frac1{[\alpha\alpha]},\frac1{[\beta\beta]},\frac1{[\gamma\gamma]}\) etc., or the mean errors to be feared in them
\[
\begin{aligned}
\text{for }x\mathord{.}\,\mathord{.}\,\mathord{.}\,\mathord{.}\,\mathord{.}\,\mathord{.}\,&m\BareRoot{p[\alpha\alpha]}=m'\BareRoot{p'[\alpha\alpha]}=m''\BareRoot{p''[\alpha\alpha]}\text{ etc.}\\
\text{for }y\mathord{.}\,\mathord{.}\,\mathord{.}\,\mathord{.}\,\mathord{.}\,\mathord{.}\,&m\BareRoot{p[\beta\beta]}=m'\BareRoot{p'[\beta\beta]}=m''\BareRoot{p''[\beta\beta]}\text{ etc.}\\
\text{for }z\mathord{.}\,\mathord{.}\,\mathord{.}\,\mathord{.}\,\mathord{.}\,\mathord{.}\,&m\BareRoot{p[\gamma\gamma]}=m'\BareRoot{p'[\gamma\gamma]}=m''\BareRoot{p''[\gamma\gamma]}\text{ etc.}
\end{aligned}
\]
\begin{center}etc.\end{center}
in agreement with what we taught in the \textit{Theoria Motus Corporum Coelestium}.

\Article{22}
It will be useful to discuss separately, in a few words, the simplest and at the same time most frequent case, in which there is only one unknown and \(V=x,V'=x,V''=x\) etc. We shall then have \(a=\BareRoot{p},a'=\BareRoot{p'},a''=\BareRoot{p''}\) etc., \(l=-L\BareRoot{p},l'=-L'\BareRoot{p'},l''=-L''\BareRoot{p''}\) etc., and hence
\SourceLeaf{031}
\[
\xi=(p+p'+p''+\text{etc.})x-(pL+p'L'+p''L''+\text{etc.})
\]
From this, further,
\[
\begin{aligned}
\relax[\alpha\alpha]&=\frac1{p+p'+p''+\text{etc.}}\\[6pt]
A&=\frac{pL+p'L'+p''L''+\text{etc.}}{p+p'+p''+\text{etc.}}
\end{aligned}
\]
If, therefore, several observations of unequal precision, whose weights are respectively \(p,p',p''\) etc., have given the value of the same quantity from the first \(=L\), from the second \(=L'\), from the third \(=L''\) etc., its most plausible value will be
\[
=\frac{pL+p'L'+p''L''+\text{etc.}}{p+p'+p''+\text{etc.}}
\]
and the weight of this determination \(=p+p'+p''\text{ etc.}\) If all observations have equal precision, the most plausible value will be
\[
=\frac{L+L'+L''+\text{etc.}}{\pi}
\]
that is, equal to the arithmetic mean of the observed values, and the weight of this determination \(=\pi\), taking the weight of the observations as the unit.
\par\nobreak\medskip\centerline{\rule{.25\linewidth}{.3pt}}
\clearpage
